\chapter{Recreations}
\label{ch:recreations}

\section{The Game of 5}

Here is a game of counting\index{Game of 5}.\footnote{This game is usually played as the ``game of $7$'' in the base-10 world, because $7$ is the most annoying one to test for divisibility. In the base-8 world, $5$ is the most annoying one.} Two or more friends can play together. Each takes a turn saying the next number. If the number is a multiple of $5$, or if it ends with $5$, clap instead of saying it. If you say it, you lose. If you ever say the nonsense word ``eight'' or ``nine,'' you lose too.

The claps in the hundreds table are circled.

\begin{center}
\setlength{\tabcolsep}{0.42em}
\renewcommand{\arraystretch}{1.55}
\begin{tabular}{|c|c|c|c|c|c|c|c|}
\hline
$1$ & $2$ & $3$ & $4$ & \sievec{5} & $6$ & $7$ & $10$ \\
\hline
$11$ & \sievec{12} & $13$ & $14$ & \sievec{15} & $16$ & \sievec{17} & $20$ \\
\hline
$21$ & $22$ & $23$ & \sievec{24} & \sievec{25} & $26$ & $27$ & $30$ \\
\hline
\sievec{31} & $32$ & $33$ & $34$ & \sievec{35} & \sievec{36} & $37$ & $40$ \\
\hline
$41$ & $42$ & \sievec{43} & $44$ & \sievec{45} & $46$ & $47$ & \sievec{50} \\
\hline
$51$ & $52$ & $53$ & $54$ & \sievec{55} & $56$ & $57$ & $60$ \\
\hline
$61$ & \sievec{62} & $63$ & $64$ & \sievec{65} & $66$ & \sievec{67} & $70$ \\
\hline
$71$ & $72$ & $73$ & \sievec{74} & \sievec{75} & $76$ & $77$ & $100$ \\
\hline
\end{tabular}
\end{center}

\section{Kaprekar's Routine}

Kaprekar's Routine\index{Kaprekar's Routine} describes a procedure to generate a sequence of numbers. You start with a three-digit number whose digits are not all the same. Rearrange the digits to make the largest number you can, and the smallest. Subtract the small one from the large one. Repeat with the result, keeping three digits (pad a leading zero if you need to).

Whatever you start with, you land on $374$, and then you stay there: $743 - 347 = 374$. For example,
\[
  157 \to 751 - 157 = 572 \to 752 - 257 = 473 \to 743 - 347 = 374.
\]

The same procedure can be carried out with four-digit numbers. Most starting numbers fall into this cycle instead of a single fixed kernel (which is less interesting, but still fascinating):
\[
  1776 \to 6062 \to 6332 \to 3774 \to 4244 \to 1776 \to \cdots
\]

\section{Narcissistic Numbers}

A narcissistic number\index{narcissistic number} is equal to the sum of its own digits, each raised to the power of how many digits there are. Single-digit numbers are narcissistic in a boring way, for example $5^{1} = 5$.

With two digits there are already real ones:
\[
  24 = 2^{2} + 4^{2}, \qquad 64 = 6^{2} + 4^{2}.
\]

The three-digit narcissistic numbers are
\begin{center}
\begin{tabular}{r@{\hspace{0.45em}}c@{\hspace{0.45em}}l}
$134$ & $=$ & $1^{3} + 3^{3} + 4^{3}$ \\
$205$ & $=$ & $2^{3} + 0^{3} + 5^{3}$ \\
$463$ & $=$ & $4^{3} + 6^{3} + 3^{3}$ \\
$660$ & $=$ & $6^{3} + 6^{3} + 0^{3}$ \\
$661$ & $=$ & $6^{3} + 6^{3} + 1^{3}$ \\
\end{tabular}
\end{center}

There are no four-digit ones. The five-digit ones are
\begin{center}
\begin{tabular}{r@{\hspace{0.45em}}c@{\hspace{0.45em}}l}
$40663$ & $=$ & $4^{5} + 0^{5} + 6^{5} + 6^{5} + 3^{5}$ \\
$42710$ & $=$ & $4^{5} + 2^{5} + 7^{5} + 1^{5} + 0^{5}$ \\
$42711$ & $=$ & $4^{5} + 2^{5} + 7^{5} + 1^{5} + 1^{5}$ \\
$60007$ & $=$ & $6^{5} + 0^{5} + 0^{5} + 0^{5} + 7^{5}$ \\
$62047$ & $=$ & $6^{5} + 2^{5} + 0^{5} + 4^{5} + 7^{5}$ \\
\end{tabular}
\end{center}

There is only one six-digit one:
\[
  636703 = 6^{6} + 3^{6} + 6^{6} + 7^{6} + 0^{6} + 3^{6}.
\]
